\documentclass[11pt, letterpaper]{article}
% ============================================================
% preamble.tex — Shared LaTeX preamble for Learning to Autolens
% ============================================================
% Contains: packages, physics notation macros, formatting.
% Include in each module via % ============================================================
% preamble.tex — Shared LaTeX preamble for Learning to Autolens
% ============================================================
% Contains: packages, physics notation macros, formatting.
% Include in each module via % ============================================================
% preamble.tex — Shared LaTeX preamble for Learning to Autolens
% ============================================================
% Contains: packages, physics notation macros, formatting.
% Include in each module via \input{../preamble}
% ============================================================

% --- Packages ---
\usepackage[utf8]{inputenc}
\usepackage[T1]{fontenc}
\usepackage{amsmath, amssymb, amsthm}
\usepackage{physics}       % \vb, \grad, \div, \curl, \dd, etc.
\usepackage{mathtools}     % \coloneqq, \prescript
\usepackage{bm}            % \bm for bold math
\usepackage{graphicx}
\usepackage{float}
\usepackage{booktabs}      % Professional tables
\usepackage{hyperref}
\usepackage{cleveref}      % \cref for smart cross-refs
\usepackage{xcolor}
\usepackage[margin=1in]{geometry}
\usepackage{enumitem}      % Custom list formatting
\usepackage{fancyhdr}      % Headers and footers
\usepackage{tcolorbox}     % Colored boxes for key results
\usepackage{listings}      % Code listings

% --- Color scheme ---
\definecolor{codeblue}{RGB}{31, 119, 180}
\definecolor{codegray}{RGB}{128, 128, 128}
\definecolor{keyresult}{RGB}{39, 123, 69}
\definecolor{exercise}{RGB}{178, 102, 0}

% --- tcolorbox environments ---
\newtcolorbox{keyresult}[1][]{
  colback=keyresult!5, colframe=keyresult!70,
  fonttitle=\bfseries, title={Key Result}, #1
}

\newtcolorbox{pythonbox}[1][]{
  colback=codeblue!5, colframe=codeblue!50,
  fonttitle=\bfseries, title={PyAutoLens Implementation}, #1
}

\newtcolorbox{exercisebox}[1][]{
  colback=exercise!5, colframe=exercise!50,
  fonttitle=\bfseries, title={Exercise}, #1
}

% --- Code listing style ---
\lstset{
  language=Python,
  basicstyle=\ttfamily\small,
  keywordstyle=\color{codeblue}\bfseries,
  commentstyle=\color{codegray}\itshape,
  stringstyle=\color{keyresult},
  breaklines=true,
  frame=single,
  framerule=0.5pt,
  rulecolor=\color{codegray!50},
}

% --- Physics macros ---
% Vectors (bold upright)
\newcommand{\vtheta}{\bm{\theta}}       % Image-plane position
\newcommand{\vbeta}{\bm{\beta}}         % Source-plane position
\newcommand{\valpha}{\bm{\alpha}}       % Deflection angle
\newcommand{\vx}{\bm{x}}               % Generic position

% Lensing quantities
\newcommand{\thetaE}{\theta_{\rm E}}    % Einstein radius
\newcommand{\SigmaCr}{\Sigma_{\rm cr}}  % Critical surface density
\newcommand{\Dd}{D_{\rm d}}             % Angular diameter distance to lens
\newcommand{\Ds}{D_{\rm s}}             % Angular diameter distance to source
\newcommand{\Dds}{D_{\rm ds}}           % Angular diameter distance lens-source

% Statistical quantities
\newcommand{\chisq}{\chi^2}             % Chi-squared
\newcommand{\chisqred}{\chi^2_{\rm red}} % Reduced chi-squared
\newcommand{\Like}{\mathcal{L}}         % Likelihood
\newcommand{\Evid}{\mathcal{Z}}         % Evidence
\newcommand{\Post}{\mathcal{P}}         % Posterior

% Abbreviations for references
\newcommand{\CK}{C\&K}                  % Congdon & Keeton (2018)
\newcommand{\NB}{N\&B}                  % Narayan & Bartelmann (1997)
\newcommand{\Night}{Nightingale+18}     % Nightingale, Dye & Massey (2018)

% --- Headers ---
\pagestyle{fancy}
\fancyhf{}
\fancyhead[L]{\small\textit{Learning to Autolens}}
\fancyhead[R]{\small\thepage}
\renewcommand{\headrulewidth}{0.4pt}

% --- Theorem environments ---
\theoremstyle{definition}
\newtheorem{definition}{Definition}[section]
\newtheorem{example}{Example}[section]

% ============================================================

% --- Packages ---
\usepackage[utf8]{inputenc}
\usepackage[T1]{fontenc}
\usepackage{amsmath, amssymb, amsthm}
\usepackage{physics}       % \vb, \grad, \div, \curl, \dd, etc.
\usepackage{mathtools}     % \coloneqq, \prescript
\usepackage{bm}            % \bm for bold math
\usepackage{graphicx}
\usepackage{float}
\usepackage{booktabs}      % Professional tables
\usepackage{hyperref}
\usepackage{cleveref}      % \cref for smart cross-refs
\usepackage{xcolor}
\usepackage[margin=1in]{geometry}
\usepackage{enumitem}      % Custom list formatting
\usepackage{fancyhdr}      % Headers and footers
\usepackage{tcolorbox}     % Colored boxes for key results
\usepackage{listings}      % Code listings

% --- Color scheme ---
\definecolor{codeblue}{RGB}{31, 119, 180}
\definecolor{codegray}{RGB}{128, 128, 128}
\definecolor{keyresult}{RGB}{39, 123, 69}
\definecolor{exercise}{RGB}{178, 102, 0}

% --- tcolorbox environments ---
\newtcolorbox{keyresult}[1][]{
  colback=keyresult!5, colframe=keyresult!70,
  fonttitle=\bfseries, title={Key Result}, #1
}

\newtcolorbox{pythonbox}[1][]{
  colback=codeblue!5, colframe=codeblue!50,
  fonttitle=\bfseries, title={PyAutoLens Implementation}, #1
}

\newtcolorbox{exercisebox}[1][]{
  colback=exercise!5, colframe=exercise!50,
  fonttitle=\bfseries, title={Exercise}, #1
}

% --- Code listing style ---
\lstset{
  language=Python,
  basicstyle=\ttfamily\small,
  keywordstyle=\color{codeblue}\bfseries,
  commentstyle=\color{codegray}\itshape,
  stringstyle=\color{keyresult},
  breaklines=true,
  frame=single,
  framerule=0.5pt,
  rulecolor=\color{codegray!50},
}

% --- Physics macros ---
% Vectors (bold upright)
\newcommand{\vtheta}{\bm{\theta}}       % Image-plane position
\newcommand{\vbeta}{\bm{\beta}}         % Source-plane position
\newcommand{\valpha}{\bm{\alpha}}       % Deflection angle
\newcommand{\vx}{\bm{x}}               % Generic position

% Lensing quantities
\newcommand{\thetaE}{\theta_{\rm E}}    % Einstein radius
\newcommand{\SigmaCr}{\Sigma_{\rm cr}}  % Critical surface density
\newcommand{\Dd}{D_{\rm d}}             % Angular diameter distance to lens
\newcommand{\Ds}{D_{\rm s}}             % Angular diameter distance to source
\newcommand{\Dds}{D_{\rm ds}}           % Angular diameter distance lens-source

% Statistical quantities
\newcommand{\chisq}{\chi^2}             % Chi-squared
\newcommand{\chisqred}{\chi^2_{\rm red}} % Reduced chi-squared
\newcommand{\Like}{\mathcal{L}}         % Likelihood
\newcommand{\Evid}{\mathcal{Z}}         % Evidence
\newcommand{\Post}{\mathcal{P}}         % Posterior

% Abbreviations for references
\newcommand{\CK}{C\&K}                  % Congdon & Keeton (2018)
\newcommand{\NB}{N\&B}                  % Narayan & Bartelmann (1997)
\newcommand{\Night}{Nightingale+18}     % Nightingale, Dye & Massey (2018)

% --- Headers ---
\pagestyle{fancy}
\fancyhf{}
\fancyhead[L]{\small\textit{Learning to Autolens}}
\fancyhead[R]{\small\thepage}
\renewcommand{\headrulewidth}{0.4pt}

% --- Theorem environments ---
\theoremstyle{definition}
\newtheorem{definition}{Definition}[section]
\newtheorem{example}{Example}[section]

% ============================================================

% --- Packages ---
\usepackage[utf8]{inputenc}
\usepackage[T1]{fontenc}
\usepackage{amsmath, amssymb, amsthm}
\usepackage{physics}       % \vb, \grad, \div, \curl, \dd, etc.
\usepackage{mathtools}     % \coloneqq, \prescript
\usepackage{bm}            % \bm for bold math
\usepackage{graphicx}
\usepackage{float}
\usepackage{booktabs}      % Professional tables
\usepackage{hyperref}
\usepackage{cleveref}      % \cref for smart cross-refs
\usepackage{xcolor}
\usepackage[margin=1in]{geometry}
\usepackage{enumitem}      % Custom list formatting
\usepackage{fancyhdr}      % Headers and footers
\usepackage{tcolorbox}     % Colored boxes for key results
\usepackage{listings}      % Code listings

% --- Color scheme ---
\definecolor{codeblue}{RGB}{31, 119, 180}
\definecolor{codegray}{RGB}{128, 128, 128}
\definecolor{keyresult}{RGB}{39, 123, 69}
\definecolor{exercise}{RGB}{178, 102, 0}

% --- tcolorbox environments ---
\newtcolorbox{keyresult}[1][]{
  colback=keyresult!5, colframe=keyresult!70,
  fonttitle=\bfseries, title={Key Result}, #1
}

\newtcolorbox{pythonbox}[1][]{
  colback=codeblue!5, colframe=codeblue!50,
  fonttitle=\bfseries, title={PyAutoLens Implementation}, #1
}

\newtcolorbox{exercisebox}[1][]{
  colback=exercise!5, colframe=exercise!50,
  fonttitle=\bfseries, title={Exercise}, #1
}

% --- Code listing style ---
\lstset{
  language=Python,
  basicstyle=\ttfamily\small,
  keywordstyle=\color{codeblue}\bfseries,
  commentstyle=\color{codegray}\itshape,
  stringstyle=\color{keyresult},
  breaklines=true,
  frame=single,
  framerule=0.5pt,
  rulecolor=\color{codegray!50},
}

% --- Physics macros ---
% Vectors (bold upright)
\newcommand{\vtheta}{\bm{\theta}}       % Image-plane position
\newcommand{\vbeta}{\bm{\beta}}         % Source-plane position
\newcommand{\valpha}{\bm{\alpha}}       % Deflection angle
\newcommand{\vx}{\bm{x}}               % Generic position

% Lensing quantities
\newcommand{\thetaE}{\theta_{\rm E}}    % Einstein radius
\newcommand{\SigmaCr}{\Sigma_{\rm cr}}  % Critical surface density
\newcommand{\Dd}{D_{\rm d}}             % Angular diameter distance to lens
\newcommand{\Ds}{D_{\rm s}}             % Angular diameter distance to source
\newcommand{\Dds}{D_{\rm ds}}           % Angular diameter distance lens-source

% Statistical quantities
\newcommand{\chisq}{\chi^2}             % Chi-squared
\newcommand{\chisqred}{\chi^2_{\rm red}} % Reduced chi-squared
\newcommand{\Like}{\mathcal{L}}         % Likelihood
\newcommand{\Evid}{\mathcal{Z}}         % Evidence
\newcommand{\Post}{\mathcal{P}}         % Posterior

% Abbreviations for references
\newcommand{\CK}{C\&K}                  % Congdon & Keeton (2018)
\newcommand{\NB}{N\&B}                  % Narayan & Bartelmann (1997)
\newcommand{\Night}{Nightingale+18}     % Nightingale, Dye & Massey (2018)

% --- Headers ---
\pagestyle{fancy}
\fancyhf{}
\fancyhead[L]{\small\textit{Learning to Autolens}}
\fancyhead[R]{\small\thepage}
\renewcommand{\headrulewidth}{0.4pt}

% --- Theorem environments ---
\theoremstyle{definition}
\newtheorem{definition}{Definition}[section]
\newtheorem{example}{Example}[section]


\fancyhead[C]{\small\textbf{Module 05: Pixelized Source Reconstructions}}

\title{\textbf{Module 05: Pixelized Source Reconstructions}\\[0.5em]
\large Theory Companion --- Learning to Autolens}
\author{Rodrigo C\'ordova Rosado\\
\small Harvard-Smithsonian Center for Astrophysics\\
\small \texttt{rodrigo.cordova\_rosado@cfa.harvard.edu}}
\date{\today}

\begin{document}
\maketitle

\begin{abstract}
This document develops the mathematical framework for non-parametric
(pixelized) source reconstructions in strong gravitational lensing.
We derive the linear inversion, regularization, and Bayesian evidence
that underpin the pixelized source fitting in PyAutoLens.
References: Suyu et al.\ (2006), Vegetti \& Koopmans (2009),
\Night\ Sec.~5.
\end{abstract}

\tableofcontents
\newpage

% ============================================================
\section{The Forward Problem as a Linear System}
\label{sec:linear-system}
% ============================================================

The observed image $\mathbf{d}$ (vector of $N_d$ image pixels) relates
to the source pixel brightnesses $\mathbf{s}$ (vector of $N_s$ source
pixels) via:

\begin{keyresult}
\begin{equation}
\mathbf{d} = \mathbf{B}\,\mathbf{F}\,\mathbf{s} + \mathbf{n}
\label{eq:linear-forward}
\end{equation}
\end{keyresult}

\noindent where:
\begin{itemize}[nosep]
  \item $\mathbf{F}$ ($N_d \times N_s$): the \textbf{lensing operator}.
    Entry $F_{ij}$ is the fractional area of image pixel $i$ that maps
    to source pixel $j$ via ray-tracing.
  \item $\mathbf{B}$ ($N_d \times N_d$): the \textbf{PSF blurring operator}.
    Implements 2D convolution with the PSF kernel.
  \item $\mathbf{n}$: noise vector with covariance
    $\mathbf{C} = \mathrm{diag}(\sigma_1^2, \ldots, \sigma_{N_d}^2)$.
\end{itemize}

We define the combined operator $\mathbf{M} = \mathbf{B}\mathbf{F}$
so that $\mathbf{d} = \mathbf{M}\mathbf{s} + \mathbf{n}$.

% ============================================================
\section{The Regularized Inversion}
\label{sec:inversion}
% ============================================================

\subsection{Unregularized Solution}

The maximum-likelihood source is obtained by minimizing $\chisq$:
\begin{equation}
\chisq(\mathbf{s}) = (\mathbf{d} - \mathbf{M}\mathbf{s})^T
\mathbf{C}^{-1} (\mathbf{d} - \mathbf{M}\mathbf{s})
\end{equation}

Setting $\partial\chisq/\partial\mathbf{s} = 0$ gives:
\begin{equation}
\hat{\mathbf{s}}_{\rm ML} = (\mathbf{M}^T \mathbf{C}^{-1} \mathbf{M})^{-1}
\mathbf{M}^T \mathbf{C}^{-1} \mathbf{d}
\label{eq:ml-solution}
\end{equation}

This is ill-conditioned when $N_s$ is large: the matrix
$\mathbf{M}^T \mathbf{C}^{-1} \mathbf{M}$ is nearly singular,
and the solution fits the noise (Suyu+06 Sec.~2.1).

\subsection{Regularization}

We add a smoothness penalty $\lambda\,\mathbf{s}^T\mathbf{H}\mathbf{s}$
to the objective function:
\begin{equation}
\text{minimize}\quad
\chisq(\mathbf{s}) + \lambda\,\mathbf{s}^T\mathbf{H}\mathbf{s}
\end{equation}

The regularized (MAP) solution is:

\begin{keyresult}
\begin{equation}
\hat{\mathbf{s}} = (\mathbf{M}^T \mathbf{C}^{-1} \mathbf{M}
+ \lambda\mathbf{H})^{-1}\,\mathbf{M}^T \mathbf{C}^{-1} \mathbf{d}
\label{eq:map-solution}
\end{equation}
\end{keyresult}

\noindent (Suyu+06 eq.~4). The term $\lambda\mathbf{H}$ regularizes
the matrix, making it well-conditioned.

\subsection{Regularization Matrix}

For \textbf{constant regularization}, $\mathbf{H}$ penalizes differences
between neighboring pixels:
\begin{equation}
H_{ij} = \begin{cases}
  N_i & \text{if } i = j \\
  -1  & \text{if pixels } i, j \text{ are neighbors} \\
  0   & \text{otherwise}
\end{cases}
\label{eq:H-constant}
\end{equation}
where $N_i$ is the number of neighbors of pixel $i$. This is the
discrete Laplacian on the mesh (Suyu+06 eq.~5).

For \textbf{adaptive brightness regularization}, $\lambda$ varies
per pixel based on the source brightness:
\begin{equation}
\lambda_i = \lambda_{\rm inner}
+ (\lambda_{\rm outer} - \lambda_{\rm inner})
\exp\!\left(-\frac{s_i^{\rm adapt}}{\sigma_s}\right)
\label{eq:adaptive-reg}
\end{equation}
where $s_i^{\rm adapt}$ is the brightness of pixel $i$ in the adapt
image (\Night\ Sec.~5.3).

% ============================================================
\section{Bayesian Evidence}
\label{sec:evidence}
% ============================================================

The marginal likelihood (evidence) for the pixelized source model,
after integrating out $\mathbf{s}$, is (Suyu+06 eq.~19):

\begin{keyresult}
\begin{equation}
\ln\Evid = -\frac{1}{2}\chisq(\hat{\mathbf{s}})
- \frac{1}{2}\hat{\mathbf{s}}^T \lambda\mathbf{H}\,\hat{\mathbf{s}}
+ \frac{1}{2}\ln\det(\lambda\mathbf{H})
- \frac{1}{2}\ln\det(\mathbf{M}^T\mathbf{C}^{-1}\mathbf{M} + \lambda\mathbf{H})
+ \text{const}
\label{eq:evidence}
\end{equation}
\end{keyresult}

The terms have intuitive meanings:
\begin{enumerate}[nosep]
  \item $-\frac{1}{2}\chisq$: reward for fitting the data
  \item $-\frac{1}{2}\hat{\mathbf{s}}^T\lambda\mathbf{H}\hat{\mathbf{s}}$:
    penalty for non-smooth source
  \item Determinant terms: Occam factor (penalizes unnecessary complexity)
\end{enumerate}

The evidence is maximized at the optimal $\lambda$ --- implementing
Occam's razor automatically.

% ============================================================
\section{Mesh Types}
\label{sec:meshes}
% ============================================================

\subsection{Rectangular}
A uniform grid in the source plane with spacing $\Delta\beta$.
Simple but wasteful: many pixels fall in empty regions.

\subsection{Delaunay Triangulation}
Given scattered points $\{\vbeta_j\}$ in the source plane, the Delaunay
triangulation maximizes the minimum angle of all triangles --- producing
well-shaped pixels that adapt to the point distribution.

\subsection{Voronoi Tessellation}
The dual of the Delaunay triangulation. Each Voronoi cell contains
all points closer to its generator than to any other. Natural pixel
shapes, but more complex to implement.

\subsection{Hilbert Image Mesh}
Places $N_s$ points along a Hilbert space-filling curve in the
\emph{image} plane, weighted by an adapt image (so more points
fall on bright arc regions). These points are then ray-traced to
the source plane, where a Delaunay triangulation is computed.
This is the default in the SLaM SOURCE PIX stage.

% ============================================================
\section{Connection to Mass Modeling}
\label{sec:mass-connection}
% ============================================================

The lensing operator $\mathbf{F}$ depends on the mass model (through
the deflection angles). When the mass model is free, the problem is
\textbf{semi-linear}: linear in $\mathbf{s}$ (solved analytically)
but non-linear in the mass parameters (searched by nested sampling).

This is the key efficiency gain: the non-linear search explores only
the mass parameters ($\sim$5--8 dimensions) while the source pixels
($\sim$500--1000 dimensions) are handled analytically.

A wrong mass model produces characteristic artifacts in the source
reconstruction: ringing, elongation, or artificial substructure.
These can be used diagnostically to assess mass model quality.

% ============================================================
\section*{References}
% ============================================================

\begin{itemize}[nosep, leftmargin=*]
  \item Dye, S.\ \& Warren, S.J.\ (2005), ApJ, 623, 31
  \item Nightingale, J.W., Dye, S.\ \& Massey, R.J.\ (2018), MNRAS, 478, 4738
  \item Suyu, S.H.\ et al.\ (2006), MNRAS, 371, 983
  \item Vegetti, S.\ \& Koopmans, L.V.E.\ (2009), MNRAS, 392, 945
\end{itemize}

\end{document}
